\documentclass[12pt,a4paper]{article}
\usepackage[utf8]{vietnam}
\usepackage{amsmath,amssymb}
\usepackage{tabularx}
\usepackage{tikz}
\usetikzlibrary{calc,angles,quotes,arrows.meta,patterns}
\usepackage{xcolor}
\usepackage{enumitem}
\usepackage{fancyhdr}
\usepackage{graphicx}
\usepackage[top=0.9cm, bottom=0.9cm, left=1.2cm, right=1.2cm]{geometry}
\usepackage[dethi]{ex_test}
\SAOPTN{widthSA=3.2,heightSA=0.62}

\setlength{\headheight}{15pt}
\setlength{\parskip}{0pt}

\pagestyle{fancy}
\fancyhf{}
\rhead{\textit{Môn Toán 9 -- Đề số 04}}
\lhead{\textit{Bộ đề ôn tập Giữa học kì I}}
\cfoot{\thepage}

\begin{document}

% --- KHUNG TIÊU ĐỀ ĐỀ THI ---
\noindent
\begin{tabularx}{\linewidth}{X r}
\textbf{TRƯỜNG THCS: \dotfill} & \textbf{ĐỀ KIỂM TRA ĐỊNH KỲ GIỮA HỌC KỲ I} \\
\textbf{Họ và tên: \dotfill} & \textbf{NĂM HỌC 2026 -- 2027} \\
\textbf{Lớp: \dotfill \quad SBD: \dotfill} & \textbf{Môn: TOÁN 9} \\
\hline
\end{tabularx}
\vspace{2mm}

\begin{center}
  \textbf{\large ĐỀ ÔN TẬP GIỮA HỌC KỲ I -- ĐỀ SỐ 04}\\
  \textit{Thời gian làm bài: 90 phút (Không kể thời gian phát đề)}
\end{center}
\vspace{1mm}

% --- PHẦN I ---
\noindent\textbf{\large PHẦN I. CÂU TRẮC NGHIỆM NHIỀU PHƯƠNG ÁN LỰA CHỌN (3,0 điểm)}

\textit{Thí sinh trả lời từ câu 1 đến câu 12. Mỗi câu hỏi thí sinh chỉ chọn một phương án duy nhất.}

\vspace{1mm}
\begin{ex}
Với điều kiện nào của $x$ thì phương trình $\frac{x+2}{x} = 3$ có nghĩa?
\choice
{$x \ge 0$}
{\True $x \neq 0$}
{$x = 0$}
{$x \le 0$}
\end{ex}

\begin{ex}
Phương trình $(x - 1)(x - 2)(x - 3) = 0$ có số nghiệm là:
\choice
{$1$}
{$2$}
{\True $3$}
{$4$}
\end{ex}

\begin{ex}
Trong các hệ thức sau, hệ thức nào không phải là phương trình bậc nhất hai ẩn?
\choice
{$3x + 2y = 0$}
{$0x + 2y = 0$}
{\True $0x + 0y = 2$}
{$x + 3y = 2$}
\end{ex}

\begin{ex}
Trong các hệ phương trình sau, hệ nào không phải là hệ hai phương trình bậc nhất hai ẩn?
\choice
{\True $\begin{cases} 2x - 3y = 0 \\ 0x + 0y = 5 \end{cases}$}
{$\begin{cases} 2x + 0y = -6 \\ 5x + 4y = -1 \end{cases}$}
{$\begin{cases} 5y = 10 \\ 3x - 2y = 6 \end{cases}$}
{$\begin{cases} 4x + 7y = -9 \\ -1{,}2x + 3y = 0 \end{cases}$}
\end{ex}

\begin{ex}
Cặp số $(1; 2)$ là nghiệm của hệ phương trình nào sau đây?
\choice
{$\begin{cases} 2x = -6 \\ 5x + 4y = -1 \end{cases}$}
{$\begin{cases} 5y = 10 \\ 3x - 2y = 6 \end{cases}$}
{\True $\begin{cases} 6y = 12 \\ 10x - 2y = 6 \end{cases}$}
{$\begin{cases} -2x + 6y = 12 \\ 10x - 2y = 6 \end{cases}$}
\end{ex}

\begin{ex}
Cho $a > b$. Trong các khẳng định sau, khẳng định nào sai?
\choice
{$a + 2 > b + 2$}
{$3a > 3b$}
{$-5a < -5b$}
{\True $-a + 2 > -b + 2$}
\end{ex}

\begin{ex}
Giá trị của biểu thức $N = \frac{\sin 58^\circ}{\cos 32^\circ} - \cos 60^\circ + \tan 37^\circ \cdot \tan 53^\circ + \sin 30^\circ$ là:
\choice
{\True $2$}
{$0$}
{$1$}
{$3$}
\end{ex}

\begin{ex}
Nghiệm của phương trình $\left(x + \frac{1}{3}\right)(x - 3) = 0$ là:
\choice
{$x = -\frac{1}{3}$}
{$x = \frac{1}{3}$}
{\True $x = -\frac{1}{3}$ và $x = 3$}
{$x = -\frac{1}{3}$ và $x = -3$}
\end{ex}

\begin{ex}
Cho tam giác $ABC$ vuông tại $A$. Hệ thức nào sau đây là đúng?
\choice
{$AB = AC \cdot \tan B$}
{$AB = AC \cdot \cos B$}
{$AB = AC \cdot \cos C$}
{\True $AB = AC \cdot \tan C$}
\end{ex}

\begin{ex}
Một ngôi nhà có bậc thềm $AH$ cao $50\text{ cm}$. Người ta làm một chiếc cầu xe $AB$ dài $1\text{ m}$. Khi đó góc tạo bởi chiếc cầu xe và mặt đất bằng:
\choice
{$60^\circ$}
{$50^\circ$}
{\True $30^\circ$}
{$45^\circ$}
\end{ex}

\begin{ex}
Nghiệm của hệ phương trình $\begin{cases} x + y = 3 \\ x - y = 1 \end{cases}$ là:
\choice
{\True $(2; 1)$}
{$(1; 2)$}
{$(3; 0)$}
{$(2; -1)$}
\end{ex}

\begin{ex}
Cho tam giác $ABC$ vuông tại $A$, đường cao $AH$. Biết $BH = 4\text{ cm}, HC = 9\text{ cm}$. Độ dài đường cao $AH$ bằng:
\choice
{\True $6\text{ cm}$}
{$36\text{ cm}$}
{$5\text{ cm}$}
{$13\text{ cm}$}
\end{ex}


\vspace{1mm}
\noindent\textbf{BẢNG TRẢ LỜI TRẮC NGHIỆM PHẦN I:}
\begin{center}
\begin{tabular}{|c|c|c|c|c|c|c|c|c|c|c|c|c|}
\hline
\textbf{Câu} & \textbf{1} & \textbf{2} & \textbf{3} & \textbf{4} & \textbf{5} & \textbf{6} & \textbf{7} & \textbf{8} & \textbf{9} & \textbf{10} & \textbf{11} & \textbf{12} \\
\hline
\textbf{Chọn} & \rule{0pt}{0.65cm} & & & & & & & & & & & \\
\hline
\end{tabular}
\end{center}
\vspace{2mm}

\setcounter{ex}{0}
\noindent\textbf{\large PHẦN II. CÂU TRẮC NGHIỆM ĐÚNG SAI (2,0 điểm)}

\textit{Thí sinh trả lời câu 1 và câu 2. Trong mỗi ý a), b), c), d) ở mỗi câu, thí sinh chọn đúng hoặc sai.}

\vspace{1mm}
\begin{ex}
Cho hệ phương trình $\begin{cases} 3x - 6y = 7 \\ 3x - 6y = 3. \end{cases}$ Xét tính Đúng/Sai của các khẳng định sau:
\choiceTF
{\True Nếu nhân hai vế của phương trình thứ hai với 2 thì hệ phương trình thay đổi giá trị nghiệm}
{\True Trừ từng vế phương trình thứ nhất cho phương trình thứ hai của hệ, ta được phương trình $0x = 4$}
{\True Phương trình $0x = 4$ là phương trình vô nghiệm}
{Hệ phương trình đã cho có nghiệm duy nhất}
\end{ex}

\begin{ex}
Cho tam giác $DEF$ vuông tại $D$, biết $DE = 6\text{ cm}$ và $DF = 8\text{ cm}$. Xét tính Đúng/Sai của các khẳng định sau:
\choiceTF
{\True Độ dài cạnh huyền $EF = 10\text{ cm}$}
{\True Tỉ số lượng giác $\tan E = \frac{4}{3}$}
{\True Số đo góc $E$ xấp xỉ bằng $53^\circ$}
{\True Độ dài đường cao xuất phát từ đỉnh $D$ bằng $4{,}8\text{ cm}$}
\end{ex}


\noindent\textbf{BẢNG TRẢ LỜI PHẦN II:}
\begin{center}
\begin{tabular}{|c|c|c|c|c||c|c|c|c|c|}
\hline
\textbf{Câu 1} & \textbf{a)} & \textbf{b)} & \textbf{c)} & \textbf{d)} & \textbf{Câu 2} & \textbf{a)} & \textbf{b)} & \textbf{c)} & \textbf{d)} \\
\hline
\textbf{Đ/S} & \rule{0pt}{0.6cm} & & & & \textbf{Đ/S} & & & & \\
\hline
\end{tabular}
\end{center}
\vspace{2mm}

\setcounter{ex}{0}
\noindent\textbf{\large PHẦN III. CÂU HỎI TRẢ LỜI NGẮN (2,0 điểm)}

\textit{Thí sinh trả lời câu 1 và câu 2. Điền kết quả vào ô trống tương ứng.}

\vspace{1mm}
\begin{ex}
Cho hệ phương trình $\begin{cases} x - 5y = 21 \\ -6x + 3y = -45. \end{cases}$ Biết $(x_0; y_0)$ là nghiệm của hệ, tính giá trị biểu thức $T = x_0 + y_0$. \shortans{$3$}
\end{ex}

\begin{ex}
Cho hai góc ở tâm kề nhau $\widehat{AOC} = 45^\circ$ và $\widehat{COB} = 30^\circ$ trên cùng một đường tròn. Số đo cung nhỏ $\wideparen{AB}$ bằng bao nhiêu độ? \shortans{$75^\circ$}
\end{ex}


\setcounter{ex}{0}
\noindent\textbf{\large PHẦN IV. TỰ LUẬN (3,0 điểm)}

\textit{Thí sinh trình bày chi tiết lời giải vào giấy làm bài.}

\vspace{1mm}
\begin{ex} (1,0 điểm)
Giải các phương trình sau:
\begin{enumerate}[label=\textbf{\alph*)}, leftmargin=0.8cm, itemsep=2pt]
\item (0,5 điểm) $(x - 3)(x + 7) = 0$
\item (0,5 điểm) $\frac{5}{x + 2} + \frac{3}{x - 1} = \frac{3x + 4}{(x + 2)(x - 1)}$
\end{enumerate}
\end{ex}
\vspace{0.5mm}

\begin{ex} (1,0 điểm)
Nhân dịp ngày Giỗ Tổ Hùng Vương, một siêu thị điện máy đã giảm giá nhiều mặt hàng để kích cầu mua sắm. Giá niêm yết của một chiếc tủ lạnh và một chiếc máy giặt có tổng số tiền là $25{,}4\text{ triệu đồng}$. Tuy nhiên, trong dịp này tủ lạnh được giảm $40\%$ giá niêm yết và máy giặt được giảm $25\%$ giá niêm yết. Vì thế, cô Liên đã mua hai mặt hàng trên với tổng số tiền là $16{,}77\text{ triệu đồng}$. Hỏi giá niêm yết của mỗi mặt hàng trên là bao nhiêu triệu đồng?
\end{ex}
\vspace{0.5mm}

\begin{ex} (1,0 điểm)
Một người đứng cách chân tháp $13{,}65\text{ m}$ nhìn lên đỉnh tháp với phương nhìn hợp với phương nằm ngang một góc bằng $58^\circ$. Biết mắt của người đó cách chân của mình một khoảng $1{,}55\text{ m}$, hỏi ngọn tháp cao bao nhiêu mét (làm tròn đến hai chữ số thập phân)?
\begin{center}
\begin{tikzpicture}[scale=1.1, line join=round, line cap=round, >=stealth]
  \coordinate (M) at (0,1.55); % Mat nguoi quan sat
  \coordinate (O) at (0,0); % Chan nguoi
  \coordinate (B) at (4.5,0); % Chan thap
  \coordinate (H) at (4.5,1.55); % Hinh chieu ngang cua mat tren thap
  \coordinate (A) at (4.5,4.7); % Dinh thap
  \draw[thick] (-0.5,0) -- (5.2,0) node[below] {Mặt đất};
  \draw[line width=1.5pt, blue!70!black] (B) -- (A) node[midway, right] {Tháp};
  \draw[thick] (O) -- (M) node[midway, left] {$1{,}55\text{ m}$};
  \draw[thick, dashed] (M) -- (H) node[midway, above] {$13{,}65\text{ m}$};
  \draw[thick, red!80!black] (M) -- (A);
  \draw (4.5,1.8) -- (4.2,1.8) -- (4.2,1.55);
  \draw (4.5,0.25) -- (4.2,0.25) -- (4.2,0);
  \draw[->, thick, purple] (1.0,1.55) arc (0:58:1);
  \node at (1.3,1.9) [purple] {$58^\circ$};
  \filldraw[black] (M) circle (1.5pt) node[left] {Mắt $(M)$};
  \filldraw[black] (A) circle (1.5pt) node[above] {Đỉnh tháp $(A)$};
  \filldraw[black] (H) circle (1.5pt) node[below right] {$H$};
  \filldraw[black] (B) circle (1.5pt) node[below right] {Chân tháp $(B)$};
\end{tikzpicture}
\end{center}
\end{ex}
\vspace{0.5mm}


\enlargethispage{1.5cm}
\begin{center}
\rule{4cm}{0.4pt}\\
\textit{--- HẾT ---}
\end{center}
\end{document}
